% Generated draft: requires mathematical and source review.
\documentclass[11pt]{article}
\usepackage[margin=25mm]{geometry}
\usepackage{fontspec}
\setmainfont{Latin Modern Roman}
\usepackage{amsmath,amssymb,mathtools,unicode-math}
\setmathfont{Latin Modern Math}
\usepackage{enumitem,needspace,xurl,hyperref}
\setlength{\parindent}{0pt}
\setlength{\parskip}{5pt}
\setlength{\emergencystretch}{3em}
\newcommand{\N}{\mathbb{N}}
\newcommand{\Z}{\mathbb{Z}}
\newcommand{\Q}{\mathbb{Q}}
\newcommand{\R}{\mathbb{R}}
\newcommand{\C}{\mathbb{C}}
\newcommand{\F}{\mathbb{F}}
\newcommand{\A}{\mathbb{A}}
\newcommand{\PP}{\mathbb{P}}
\newcommand{\E}{\mathbb{E}}
\newcommand{\eps}{\varepsilon}
\newcommand{\dd}{\,\mathrm{d}}
\begin{document}
\section*{UnsolvedMath ID: 701147}
\textbf{Source problem code:} catalogue entry
\subsection{Definitions and mathematical statement}
A profinite group is a topological group isomorphic to an inverse limit of finite discrete groups. A free profinite basis of a profinite group \(F\) is a subset \(X\subseteq F\) which topologically generates \(F\), converges to the identity (every identity neighbourhood contains all but finitely many elements of \(X\)), and has the following universal property: for every profinite group \(H\), every map \(f:X\to H\) converging to the identity extends uniquely to a continuous homomorphism \(F\to H\). A profinite group is free if it has such a basis. It is free of finite rank \(r\) if it has such a basis with \(r\) elements, where \(r\in\Z_{\geq 0}\); rank zero gives the trivial group.

\textbf{Existence question (O. V. Mel'nikov).} Does there exist a profinite group \(G\) which is not a free profinite group of any rank, together with a nonempty family \((N_\alpha)_{\alpha\in I}\) of closed normal subgroups of \(G\), directed by reverse inclusion, such that the canonical homomorphism is an isomorphism of topological groups
\[
G\longrightarrow\varprojlim_{\alpha\in I}G/N_\alpha,
\qquad g\longmapsto(gN_\alpha)_{\alpha\in I},
\]
and every quotient \(G/N_\alpha\), with its quotient topology, is a free profinite group of finite rank? Here directedness means that for every \(\alpha,\beta\in I\), some \(\gamma\in I\) satisfies \(N_\gamma\subseteq N_\alpha\cap N_\beta\); the bonding homomorphisms are the natural quotient maps. The finite ranks may depend on \(\alpha\); no uniform bound is stipulated.

\textbf{Supplied qualifications.} The source states that the finiteness restriction on these ranks is essential. It also states that any such nonfree group \(G\) cannot satisfy the first axiom of countability, citing Mel'nikov's 1980 paper listed below. First countability here means that the identity has a countable neighbourhood basis, equivalently that every point does. These are preserved source assertions, not independently verified conclusions of this draft.
\subsection{Short English statement}
Does a nonfree profinite group exist which is canonically an inverse limit of its finite-rank free profinite quotients? The ranks may vary. The supplied source says that finite rank is essential and that any example must fail first countability.
\subsection{Sources}
\begin{itemize}
\item Supplied local UnsolvedMath registry record: permanent ID \(701147\), exact source code \texttt{catalogue entry}, title \emph{Kourovka Notebook Problem 15.59}. The registry labels the problem \emph{open}; that label is supplied metadata, not an independently verified mathematical or literature assessment.
\item \emph{The Kourovka Notebook (Unsolved Problems in Group Theory)}, original Problem 15.59, 15th issue (2002), proposed by O. V. Mel'nikov. The supplied record identifies the later source as the 21st issue edition (2026), \texttt{arXiv:1401.0300v48}, and identifies E. I. Khukhro and V. D. Mazurov as editors of the collection.
\item O. V. Mel'nikov, \emph{Dokl. AN BSSR}, \textbf{24}, no. 11 (1980), 968--970, in Russian. This reference is supplied in connection with the first-countability obstruction; no article title is supplied.
\end{itemize}
\end{document}
